\section{Discussion}
\label{sec:discussion}
We discuss the trade-offs of the overlay approach, portability to other Cray/HPE
systems, and limitations observed during validation.

\subsection{Limitations}
\label{sec:limitations}
The overlay depends on PE wrappers and module hierarchy conventions; changes in
those interfaces require updates to Afar-prgenv. Some vendor libraries ship
Fortran modules that are not compatible with AFAR compilers and require shims
or local builds. Accelerator libraries (e.g., `libsci-acc`) may also lag when
runtime dependencies are tied to older ROCm versions.
Parallel HDF5 validation can also be constrained by missing C++ `pkg-config`
metadata in the PE environment, which must be addressed before full coverage.

We also intentionally avoid performance benchmarking in this work. Our goal is
functional readiness and compatibility; performance evaluation is deferred
until after a drop passes the compatibility gate.

\subsection{Portability}
\label{sec:portability}
Afar-prgenv generalizes to other Cray/HPE systems because it leverages the same
PE module hierarchy and wrappers. The system-specific configuration is limited
to the AFAR root path, version mappings, and the Cray MPICH install path, all
of which are captured in configuration files.

\subsection{Lessons Learned}
\label{sec:lessons}
Two operational lessons stand out. First, early drop readiness is mostly a
metadata problem: wrapper semantics and `pkg-config` paths must be correct
before any application testing can succeed. Second, reproducible repro tests
are critical for support teams because they isolate regressions to specific
libraries or wrapper behaviors. These lessons inform how we prioritize future
automation and validation.
